\section{Related Work}
Predictive maintenance and prognostics have been studied extensively in the context of run-to-failure datasets, especially the C-MAPSS turbofan benchmarks introduced by NASA for degradation modeling and remaining useful life prediction \cite{Saxena2008CMAPSS,Ramasso2014CMAPSS}. These datasets remain popular because they combine multivariate sensor streams, multiple fault modes, and a clear supervised evaluation target.

The IMS bearing data provide a complementary failure setting that is more directly tied to vibration analysis and bearing health monitoring \cite{NASAIMSBearings}. Unlike C-MAPSS, IMS does not provide ready-made labels, so it is often used to study anomaly detection, fault progression, and failure-stage estimation. That makes it useful here as a secondary benchmark for testing whether a model can generalize beyond a single canonical prognostics dataset.

On the modeling side, this project is motivated by adaptive computation methods that allocate additional internal processing only when the input is difficult. Adaptive Computation Time introduced learned halting for recurrent networks \cite{Graves2016ACT}, while Universal Transformers extended recurrence with repeated refinement of latent states \cite{Dehghani2018UniversalTransformers}. More recent looped reasoning work argues that iterative latent computation can improve reasoning quality by using depth as a learned resource rather than a fixed architectural constant \cite{Zhu2025Ouro, Banino2021PonderNet}.

The present project connects these strands by asking whether adaptive looped reasoning can improve predictive maintenance decisions. The key question is not simply whether the model can classify healthy versus failing equipment, but whether it can do so with a useful compute profile that spends more effort on ambiguous windows and less effort on obvious ones.
